\documentclass[10pt,a4paper]{amsart}
\usepackage[margin=19mm]{geometry}
\usepackage{amsmath,amssymb,mathtools}
\usepackage[scr=rsfs,cal=euler]{mathalfa}
\usepackage{xfrac}
\usepackage[unicode,hidelinks,bookmarks=false]{hyperref}
\newcommand{\ie}{i.e.\ }
\newcommand{\cf}{cf.\ }
\newcommand{\resp}{respectively\ }
\newcommand{\Subsection}{\subsection}
\providecommand{\nobreakdash}{\nobreak\hbox{-}}
\setlength{\emergencystretch}{2em}
\multlinegap0pt
\usepackage{bm}
\usepackage{amssymb}
\newtheorem{theorem}{Theorem}
\newtheorem{corollary}{Corollary}
\title{Product systems arising from L\'evy processes}
\author{Remus Floricel, Peter Wadel}
\date{Source: arXiv:2410.14893v2 (CC BY 4.0)}
\begin{document}
\maketitle

\noindent\emph{Source and licence.} Product systems arising from L\'evy processes. Version arXiv:2410.14893v2 (CC BY 4.0), distributed by arXiv under the Creative Commons Attribution 4.0 International licence, http://creativecommons.org/licenses/by/4.0/, as recorded in the arXiv licence metadata for this exact version and shown on the abstract page https://arxiv.org/abs/2410.14893v2. Full licence notice: \texttt{licenses/arxiv-2410.14893-CC-BY-4.0.txt}.\par\smallskip

\noindent\textit{Editorial scope.} Counted unit: the article's theorem theorem:completely\_spatial (source0.tex lines 590--592). The article's complete original proof of it is delivered with it (source0.tex lines 594--700, one \texttt{proof} environment of 107 lines). No mathematics is written, repaired or re-derived: the statement and the proof are the article's own lines, in the article's own notation, and no retained line is substituted or re-flowed.  Retained uncounted as the source-local context the proof needs: lines 509--511; lines 574--576.  Nothing was omitted from any retained span, so the package carries no omission record.  No retained line contains a citation, so the wrapper carries no bibliography and no bibliography entry.  No retained line contains a figure environment, an image-inclusion command or drawing code, so no image is delivered and the package declares no asset dependency.  Editorial formatting only, in addition: an AMS \texttt{amsart} preamble that restates the article's own declarations and macros used by the retained lines, namely the environments \texttt{theorem}, \texttt{corollary} on separate counters, together with the standard packages those lines need.  The retained environments print in this document as Theorem 1, Corollary 1 in order of appearance, whereas the article prints the same items with the numbering of its own full document.  The complete native source of the article as distributed in the arXiv e-print for this exact version is bundled unchanged as \texttt{sources/167087/source0.tex}.
\begin{theorem}\label{main1}
Let \( \bm{L} = \{L_t\}_{t > 0} \) be a purely Gaussian L\'evy process in a real, separable Banach space \( B \). Then any unit of \( E^{\bm{L}} \) is equivalent to a standard exponential unit.
\end{theorem}

\begin{corollary}\label{cor:extension}
Let \( \bm{L} = \{L_t\}_{t > 0} \) be a Gaussian L\'evy process in a real, separable Banach space \( B \), with drift \( b \in B \), non-degenerate covariance \( Q \), and zero L\'evy measure. Then any unit of the product system \( E^{\bm{L}} \) is equivalent to a standard exponential unit. 
\end{corollary}

\begin{theorem}\label{theorem:completely_spatial}
Let \( \bm{L} = \{ L_t \}_{t > 0} \) be a Gaussian L\'evy process in a real, separable Banach space \( B \), with drift \( b \in B \), non-degenerate covariance \( Q \), and zero L\'evy measure. Then the associated product system \( E^{\bm{L}} \) is completely spatial.
\end{theorem}

\begin{proof} First, we notice that the presence of a drift \( b \) does not affect the structure of the product system. Indeed, as in the proof of Corollary \ref{cor:extension}, consider the process \( \bm{L}^0\) obtained by removing the drift from  \( \bm{L}\). Then the Hilbert spaces \( E^{\bm{L}}(t)\)  and \( E^{\bm{L}^0}(t)\)  coincide because the shift by \( t b \) is deterministic and does not alter the \( L^2 \) structure. Thus, without loss of generality, we may assume \( b = 0 \)  in the remainder of the proof.

To prove that the product system \( E^{\bm{L}} \) is completely spatial, we must show that for each \( t > 0 \), the Hilbert space \( E(t) = L^2(\Omega, \mathcal{F}_t, \mathbb{P}) \) is densely spanned by finite products of units. Specifically, we aim to prove that the linear span of elements of the form
\(
u^{h_1}(t_1) u^{h_2}(t_2) \dots u^{h_n}(t_n),
\)
where \( t_1 + t_2 + \dots + t_n = t \) and \( h_i \in H \), is dense in \( E(t) \).
Recall from Theorem \ref{main1} that the units of \( E^{\bm{L}} \) are given by the standard exponential units \(u^h\), \(h\in H\). Under the unitary \(
L^2(B, \mu_t)\ni f  \mapsto f \circ L_t\in E(t)\), these units correspond to the functions:
\[
\phi_h(x) = \exp\left( \langle \hat{h}, x \rangle - \frac{t}{2} \| h \|_H^2 \right), \quad x \in B.
\]

Since cylindrical functions are dense in \( L^2(B, \mu_t) \), it suffices to approximate any cylindrical function using finite linear combinations of products of the exponential functions \( \phi_h(x) \). For this, let \( f \) be a cylindrical function on \( B \), meaning there exist \( n \in \mathbb{N} \), continuous linear functionals \( \psi_1, \psi_2, \dots, \psi_n \in B^* \), and a function \( F: \mathbb{R}^n \to \mathbb{C} \) such that
\[
f(x) = F\big( \psi_1(x), \psi_2(x), \dots, \psi_n(x) \big), \quad \forall x \in B.
\]
Assume that \( F \in L^2(\mathbb{R}^n, \gamma_n) \), where \( \gamma_n \) is the standard Gaussian measure on \( \mathbb{R}^n \) with mean zero and covariance matrix \( t \Sigma \), and \( \Sigma=(\Sigma_{ij})_{1\leq i,j\leq n} \) is the covariance matrix determined by \( \psi_i \) and the covariance \( Q \), i.e., \(
\Sigma_{ij} = Q(\psi_i, \psi_j).
\) Since the Hermite polynomials \( \{ H_\alpha \} \) form an orthogonal basis in \( L^2(\mathbb{R}^n, \gamma_n) \), the function \( F \) can be expanded as
\[
F(y) = \sum_{\alpha \in \mathbb{N}_0^n} c_\alpha H_\alpha(y),
\]
where \( \alpha = (\alpha_1, \dots, \alpha_n) \) is a multi-index, \( c_\alpha \in \mathbb{C} \), and \( H_\alpha(y) = \prod_{i=1}^n H_{\alpha_i}(y_i) \), for every \(y=(y_1,\cdots, y_n)\in\mathbb{R}^n.\)   Consequently,
\[
f(x) = \sum_{\alpha \in \mathbb{N}_0^n} c_\alpha H_\alpha\big( \psi_1(x), \dots, \psi_n(x) \big).
\]
Recall that the multivariate Hermite polynomials \( H_\alpha(y) \) have the following standard generating function:
\[
\exp\left( \langle z, y \rangle - \tfrac{1}{2} \| z \|^2 \right) = \sum_{\alpha \in \mathbb{N}_0^n} \frac{z^\alpha}{\alpha!} H_\alpha(y),
\]
where \( z \in \mathbb{R}^n \), \( y \in \mathbb{R}^n \), \( z^\alpha = \prod_{i=1}^n z_i^{\alpha_i} \), and \( \alpha! = \prod_{i=1}^n \alpha_i! \).
We claim that any multivariate Hermite polynomial \( H_\alpha(y) \) can be expressed as a finite linear combination of products of exponential functions \( \phi_h(x) \). For this purpose, for each \( z \in \mathbb{R}^n \), define \( h_z \in H \) such that
\[
\hat{h}_z = \sum_{i=1}^n z_i \psi_i.
\]
Then
\begin{eqnarray*}
\phi_{h_z}(x) &=& \exp\left( \langle \hat{h}_z, x \rangle - \tfrac{t}{2} \| h_z \|_H^2 \right) = \exp\left( \langle z, y \rangle - \tfrac{t}{2} \| h_z \|_H^2 \right)
\\&=& \exp\left(  z^\top y - \tfrac{t}{2} z^\top \Sigma z \right),\end{eqnarray*} where \( y = (\psi_1(x), \dots, \psi_n(x)) \).  To align with the standard generating function of Hermite polynomials (which assumes identity covariance), we perform a change of variables to normalize the covariance matrix. By defining
\[
\tilde{y} = \frac{1}{\sqrt{t}} \Sigma^{-1/2} y, \quad \tilde{z} = \sqrt{t} \Sigma^{1/2} z
\]
we then have \[
\phi_{h_z}(x) = \exp\left( \langle \tilde{z}, \tilde{y} \rangle - \tfrac{1}{2} \| \tilde{z} \|^2 \right) = \sum_{\alpha \in \mathbb{N}_0^n} \frac{\tilde{z}^\alpha}{\alpha!} H_\alpha(\tilde{y}).
\]
This expansion allows us to express Hermite polynomials in terms of exponential functions \( \phi_h(x) \). To express \( H_\alpha(\tilde{y}) \) in terms of \( \phi_h(x) \), we consider a finite set of vectors \( \tilde{z}_1, \dots, \tilde{z}_N \in \mathbb{R}^n \) and set up a linear system. For each \( \tilde{z}_k \), we have
\[
\phi_{h_{z_k}}(x) = \sum_{\beta \in \mathbb{N}_0^n} \frac{\tilde{z}_k^\beta}{\beta!} H_\beta(\tilde{y}).
\]
Collecting these equations for \( k = 1, \dots, N \), we obtain a linear system:
\[
\begin{pmatrix}
\phi_{h_{z_1}}(x) \\
\phi_{h_{z_2}}(x) \\
\vdots \\
\phi_{h_{z_N}}(x)
\end{pmatrix}
=
A
\begin{pmatrix}
H_{\beta_1}(\tilde{y}) \\
H_{\beta_2}(\tilde{y}) \\
\vdots \\
H_{\beta_M}(\tilde{y})
\end{pmatrix},
\]
where \( A \) is the \( N \times M \) matrix with entries:
\(
A_{k, \beta} = \frac{\tilde{z}_k^\beta}{\beta!}.
\) Here, \( M \) is the number of multi-indices \( \beta \) up to degree \( |\alpha| \).

If we choose \( N = M \), i.e., the number of exponential functions equals the number of Hermite polynomials up to degree \( |\alpha| \), and ensure that the matrix \( A \) is invertible, we can solve for \( H_\alpha(\tilde{y}) \):
\[
\begin{pmatrix}
H_{\beta_1}(\tilde{y}) \\
H_{\beta_2}(\tilde{y}) \\
\vdots \\
H_{\beta_M}(\tilde{y})
\end{pmatrix}
=
A^{-1}
\begin{pmatrix}
\phi_{h_{z_1}}(x) \\
\phi_{h_{z_2}}(x) \\
\vdots \\
\phi_{h_{z_N}}(x)
\end{pmatrix}.
\]
This expresses each Hermite polynomial \( H_\alpha(\tilde{y}) \) as a finite linear combination of exponential functions \( \phi_{h_{z_k}}(x) \).

To ensure that \( A \) is invertible, we need to choose \( \tilde{z}_k \) such that the matrix \( A \) has full rank. One effective strategy is to select \( \tilde{z}_k \) as distinct standard basis vectors and their combinations. For example, for \( n = 2 \) and degrees up to \( |\alpha| = 2 \), we might choose:
\[
\tilde{z}_1 = (1, 0), \quad \tilde{z}_2 = (0, 1), \quad \tilde{z}_3 = (1, 1).
\]
With these choices, the matrix \( A \) will be invertible, allowing us to solve the linear system.


Returning to the cylindrical function \( f(x) \), we can now express it as:
\[
f(x) = \sum_{\alpha \in \mathbb{N}_0^n} c_\alpha H_\alpha\left( \tilde{y} \right) = \sum_{\alpha \in \mathbb{N}_0^n} c_\alpha \sum_{k=1}^N \left( A^{-1} \right)_{\alpha, k} \phi_{h_{z_k}}(x).
\]
Thus, \( f(x) \) is expressed as a finite linear combination of exponential functions \( \phi_{h_{z_k}}(x) \).


Since any cylindrical function \( f \) can be approximated in \( L^2(B, \mu_t) \) by finite linear combinations of exponential functions \( \phi_h(x) \), and cylindrical functions are dense in \( L^2(B, \mu_t) \), it follows that the product system \( E^{\bm{L}} \) is completely spatial.
\end{proof}

\end{document}
